\section{Примеры работы программы}
\label{sec:examples}

\subsection{Запуск программы}

После запуска из каталога проекта программа выводит консольное меню. В нём доступны запуск анализа, настройка параметров, просмотр текущей конфигурации и выход из программы. Главное меню показано на рисунке~\ref{fig:img/main-menu.png}.

\screenshot[0.75]{img/main-menu.png}{Главное меню программы}

\subsection{Настройка параметров анализа}

Пункт меню \texttt{2} используется для выбора метрик, которые будут учитываться при поиске аномалий. Пример настройки параметров анализа приведён на рисунке~\ref{fig:img/analysis-options.png}.

\screenshot[0.75]{img/analysis-options.png}{Настройка параметров анализа}

Текущие значения параметров можно вывести через пункт меню \texttt{3}. Результат вывода текущей конфигурации показан на рисунке~\ref{fig:img/current-options.png}.

\screenshot[0.65]{img/current-options.png}{Просмотр текущих параметров анализа}

Если пользователь вводит некорректный ответ при настройке, программа сообщает об ошибке и повторяет вопрос. Такой сценарий показан на рисунке~\ref{fig:img/invalid-option-input.png}.

\screenshot[0.75]{img/invalid-option-input.png}{Обработка некорректного ответа при настройке}

\subsection{Запуск анализа BDF-файла}

Для запуска анализа пользователь выбирает пункт меню \texttt{1}, вводит путь к анализируемому BDF-файлу и путь для сохранения отчёта. Если путь отчёта оставить пустым, используется путь по умолчанию. Пример запуска анализа показан на рисунке~\ref{fig:img/analysis-run.png}.

\screenshot[0.75]{img/analysis-run.png}{Запуск анализа BDF-файла}

При ошибке чтения файла программа выводит сообщение и возвращает пользователя в основной сценарий. Пример ошибки для несуществующего файла приведён на рисунке~\ref{fig:img/file-read-error.png}.

\screenshot[0.75]{img/file-read-error.png}{Ошибка чтения файла}

Если файл существует, но не является корректным BDF-файлом, программа выводит ошибку разбора. Пример такой ошибки показан на рисунке~\ref{fig:img/invalid-bdf-error.png}.

\screenshot[0.75]{img/invalid-bdf-error.png}{Ошибка разбора BDF-файла}

\subsection{Итоговый отчёт}

После успешного анализа программа сохраняет текстовый отчёт. В начале отчёта выводится общая информация о шрифте, количество глифов, число найденных аномалий и количество предложенных замен. Начало отчёта показано на рисунке~\ref{fig:img/report-summary.png}.

\screenshot[1]{img/report-summary.png}{Начало сформированного отчёта}

Далее в отчёте приводится таблица метрик по глифам. Для каждого глифа дополнительно выводится его пиксельное представление, а при нарушении различимости указывается визуально похожий глиф. Такой фрагмент отчёта приведён на рисунке~\ref{fig:img/report-glyph-metrics.png}.

\screenshot[1]{img/report-glyph-metrics.png}{Фрагмент таблицы метрик глифов}

Отдельный раздел отчёта содержит список найденных аномалий и критерии, которые были нарушены. Пример раздела с аномалиями показан на рисунке~\ref{fig:img/report-anomalies.png}.

\screenshot[1]{img/report-anomalies.png}{Список найденных аномалий}

В конце отчёта выводятся предложения замены из эталонного шрифта. Для сравнения программа показывает исходный глиф и соответствующий глиф Terminus. Пример раздела замен приведён на рисунке~\ref{fig:img/report-replacements.png}.

\screenshot[1]{img/report-replacements.png}{Предложения замены с изображениями глифов}

\newpage
